\newpage
\chapter{Calibracion del software}
\vspace{3\baselineskip}

\section{Validación y Reportes (report.py y Testing)}
El proceso de aseguramiento de la calidad del software se ejecuta mediante una suite de pruebas automatizadas que inyectan las formas de onda sintéticas y grabadas definidas en la norma IEC 61083-2 \cite{iec61083_2}. El objetivo de este procedimiento es cuantificar el sesgo del algoritmo respecto a las magnitudes de referencia metrológica.

\subsection{Flujo de Validación del Software contra TDG}
La validación implementa el protocolo de ensayos de desempeño (performance tests) estipulado en la sección 6.2 de la norma. Se cargan los casos de prueba del generador de datos de prueba (TDG) correspondientes al grupo Full Lightning Impulse (LI-A1 a LI-A12 y LI-M1 a LI-M17) e impulsos cortados en la cola (LIC-M4 y LIC-M5). El sistema procesa los archivos ASCII de dos columnas en segundos y voltios, ejecuta el pipeline algorítmico de filtrado y remoción de componentes continuas, y contrasta los resultados con los límites normativos.

\subsection{Estimación de Incertidumbre del Software (Anexo B.3.3)}
En lugar de adoptar los límites simplificados fijos, el módulo de pruebas calcula las incertidumbres combinadas del software basándose en las desviaciones analíticas reales encontradas frente al TDG, siguiendo de forma estricta las directrices matemáticas del Anexo B de la norma IEC 61083-2 \cite{iec61083_2}.

La primera componente del cálculo representa la incertidumbre estándar tipo B del software ($u_{B71}$), calculada a partir de la máxima desviación relativa observada en el conjunto de datos analizados ($n$ ondas evaluadas):

\begin{equation}
u_{B71} = \frac{1}{\sqrt{3}} \max_{i=1}^{n} \left| \frac{x_i - x_{\text{REF},i}}{x_{\text{REF},i}} \right|
\label{eq:u_B71}
\end{equation}

La segunda componente ($u_{B72}$) considera la incertidumbre expandida intrínseca del vector de referencia del TDG suministrada por los comités internacionales de mantenimiento, calibrada a un nivel de confianza del $\SI{95}{\percent}$ con un factor de cobertura $k=2$:

\begin{equation}
u_{B72} = \frac{1}{2} \max_{i=1}^{n} U_{x,i}
\label{eq:u_B72}
\end{equation}

La incertidumbre estándar combinada asignada al componente del software ($u_{B7}$) se obtiene mediante la raíz cuadrada de la suma de los cuadrados de ambas componentes independientes:

\begin{equation}
u_{B7} = \sqrt{u_{B71}^2 + u_{B72}^2}
\label{eq:u_B7_combinada}
\end{equation}

\subsection{Validaciones Paramétricas bajo IEC 61083-2}
El software ejecuta la validación de la función de transferencia de los algoritmos mediante el procesamiento de vectores digitales de referencia. El objetivo de este procedimiento DSP es aislar la curva base del impulso de cualquier perturbación de alta frecuencia (\textit{overshoot} o ruidos de conmutación) para construir la curva de tensión de ensayo de manera reproducible. La norma define diferencias máximas aceptables entre los valores verdaderos del generador de datos de prueba (TDG) y los calculados por el algoritmo bajo evaluación. El procesamiento se divide de acuerdo a la topología de la señal bajo análisis:

\begin{table}[H]
    \centering
    \caption{Tolerancias de aceptación paramétricas según la norma IEC 61083-2.}
    \label{tab:tolerancias_iec61083}
    \begin{tabular}{@{}lll@{}}
        \toprule
        \textbf{Tipo de Impulso} & \textbf{Parámetro Físico} & \textbf{Tolerancia Máxima}           \\ \midrule
        Completo (LI)            & Valor Pico ($U_t$)        & $\pm\SI{0,10}{\percent}$             \\ \cmidrule(l){2-3} 
                                 & Tiempo de Frente ($T_1$)  & $\pm\SI{2,0}{\percent}$              \\ \cmidrule(l){2-3} 
                                 & Tiempo de Cola ($T_2$)    & $\pm\SI{1,0}{\percent}$              \\ \cmidrule(l){2-3} 
                                 & Sobrepasamiento ($\beta$) & $\pm\SI{1,0}{\percent}$ (Absoluto)   \\ \midrule
        Cortado (LIC)            & Tiempo de Corte ($T_c$)   & $\pm\SI{2,0}{\percent}$              \\ \bottomrule
    \end{tabular}
\end{table}